\subsection{Point sections and generalized-Jacobian Kummer classes}

We use left Galois cocycles: $z(\sigma\tau)=z(\sigma)+\sigma z(\tau)$.
For a geometric path $\gamma$ from a rational base point $b$ to a
rational point $x$, transport back to $b$ defines the section $s_x^\gamma$.
Our section-difference convention is $s_x^\gamma(\sigma)s_b(\sigma)^{-1}$.
The geometric abelian quotient is identified with torsion of the Jacobian
by the monodromy of multiplication covers. These choices fix the sign.

\begin{claim}{B.68}{Proved}{The point-to-section Kummer comparison, including a modulus}
Let $K$ be any field, let $X/K$ be a smooth, proper, geometrically
connected curve, and let $\mathfrak m$ be an effective divisor on $X$.
Put $U=X-\operatorname{supp}(\mathfrak m)$ and assume $b\in U(K)$.
Let $J_{\mathfrak m}$ be the generalized Jacobian, with $J_0=J_X$, and
let $j_b:U\to J_{\mathfrak m}$ send $x$ to the modulus class $[x-b]$.
For every positive integer $n$ invertible in $K$, the sequence
\[
 0\longrightarrow J_{\mathfrak m}[n]\longrightarrow
 J_{\mathfrak m}\xrightarrow{[n]}J_{\mathfrak m}\longrightarrow0
\]
is exact as a sequence of fppf sheaves, and $[n]$ is finite \emph{\'etale}.
The Abel map induces a $G_K$-equivariant isomorphism of finite
abelian groups
\[
 \pi_1(U_{\overline K},\bar b)^{\mathrm{ab}}/n
 \ \xrightarrow{\sim}\ J_{\mathfrak m}[n](\overline K).
\]
For every $x\in U(K)$ the abelianized section-difference class is
\[
 [s_x-s_b]_n=\delta_n(j_b(x))\quad\hbox{in}\quad
 H^1(G_K,J_{\mathfrak m}[n](\overline K)).                 \tag{N1.1}
\]
Thus the sign is positive for $x-b$ and the stated difference convention.
Changing the geometric path changes the cocycle by a coboundary;
conjugacy throughout is conjugacy by the geometric kernel. Reversing
the section difference gives $-\delta_n(j_b(x))$.
\end{claim}
\begin{proof}
We give the geometric comparison and the arithmetic sign separately.
Over an algebraically closed field a generalized Jacobian is an
extension of the ordinary Jacobian by the connected affine group
of units of the modulus, modulo constant units. For reduced modulus
with $d>0$ geometric points its affine part is $\mathbf G_m^{d-1}$.
Increasing the multiplicities adds a connected unipotent kernel.
Multiplication by $n$ is an automorphism on that unipotent kernel:
in characteristic zero one uses a filtration by additive groups;
in characteristic $p$ one uses the finite $p$-power exponent
filtration of the units $1+\mathfrak a$ of a finite local algebra.
For the torus and the abelian quotient multiplication by $n$ is
finite, surjective and \emph{\'etale}. The same follows for
$J_{\mathfrak m}$ by descent and the extension sequence. In
particular all its geometric $n$-division points are separable.

Here is a direct verification that the indicated abelian quotient
is the one represented by these covers. For a nonempty reduced
geometric boundary $D$ the localization sequence gives
\[
 0\to H^1(X_{\overline K},\mathbf Z/n)\to
 H^1(U_{\overline K},\mathbf Z/n)\to
 (\mathbf Z/n)(-1)^D\xrightarrow{\sum}
 (\mathbf Z/n)(-1)\to0.                                \tag{N1.2}
\]
The residue map is the tame-inertia character at each boundary point;
the last map is degree. Its exactness follows from the Kummer
localization sequence, purity for a smooth curve, and the degree
isomorphism in degree two. The proper Abel map identifies the first
term with the character group of $J_X[n]$: this is
\citet[Proposition 9.1]{MilneJac}, stated there for a complete smooth
curve over a separably closed field, with all connected abelian
\emph{\'etale} covers included. On the generalized-Jacobian side the
exact sequence of $n$-torsion is
\[
 0\to (\boldsymbol\mu_n^D/\boldsymbol\mu_n)
 \to J_D[n]\to J_X[n]\to0.
\]
Its character sequence is (N1.2). To check the comparison at the
boundary, trivialize a modulus line bundle near the boundary and
move its zero once around a local parameter. The trivialization
changes in precisely that boundary factor of $\boldsymbol\mu_n^D$;
simultaneous change of all trivializations is a constant unit and
is quotiented out. Equivalently the local multiplication cover is
$w^n=z$, with residue $+1$ in the corresponding factor. Thus the
residue maps agree, and the two end isomorphisms give the middle
isomorphism. If $D$ is empty this is just the proper assertion.
The unipotent kernel contributes no $n$-torsion and no finite
abelian cover of exponent dividing $n$, so the answer is unchanged
by multiplicities. All constructions are canonical under field
automorphisms and hence $G_K$-equivariant. For an imperfect ground
field this verification can be made after its perfect closure;
purely inseparable extension induces equivalences on finite
\emph{\'etale} schemes and identifies absolute Galois groups, so
the same assertion descends to $K$.

Put $a=j_b(x)$ and choose $t\in J_{\mathfrak m}(K^{\mathrm{sep}})$
with $nt=a$. By definition of the connecting map,
\[
               \delta_n(a)(\sigma)=\sigma(t)-t.           \tag{N1.3}
\]
Pull back $[n]$ by $j_b$ to a finite \emph{\'etale}
$J_{\mathfrak m}[n]$-torsor on $U$. Its fibre at $b$ is pointed
by $0$ and its fibre at $x$ contains $t$. Choose a geometric path
$\gamma$ carrying $0$ to $t$. Such a path exists because a
geometric path identifies the fibre functors, and its possible
values on this torsor are all its geometric fibre points by the
just-proved geometric quotient isomorphism. The conjugate path
$\sigma\gamma$ carries $0$ to $\sigma(t)$. Consequently
$\gamma^{-1}(\sigma\gamma)$, acting on the fibre at $b$, has
translation coordinate $\sigma(t)-t$. With transport of sections
back to $b$, this path difference is exactly the image of
$s_x^\gamma(\sigma)s_b(\sigma)^{-1}$. This proves (N1.1).
Replacing $t$ by $t+u$ adds $\sigma(u)-u$, which also verifies
explicitly the ambiguity under a different path. Restriction to
any field extension, maps of pointed curves, and transition maps
$n'\to n$ commute with this calculation. Hence it proves the
compatible Tate-module statement as well as each finite level.
\end{proof}

\paragraph{Check record.}
Checked only in this writing: GPT-6 Astra (OpenAI, ChatGPT Pro mode),
18 November 2026, independent cocycle derivation and direct checking
of the generalized-Jacobian comparison. The source's proper-curve
Proposition 9.1 and its generalized-Jacobian discussion, including
Proposition 9.7, were opened. A separate reasoning-effort setting
was not exposed. No checking by another model family is claimed.

\paragraph{Manuscript scope and the Kummer kernels.}
Cote 43, scan 113 (manuscript page 112), contains the local
decomposition and inertia homomorphisms in a deleted passage; it
does not specify the path or cocycle sign in (N1.1). Scans 132--133
(pages 131--132) contain the Kummer sequence and the $T(J)$ bridge.
These are compatible with (N1.1), but we do not attribute our
modern sign convention to an undeciphered symbol. Cote 149,
scan 33 (manuscript page 16), explicitly identifies the
point homomorphism with the inverse-limit Kummer homomorphism.
Scans 34--36 then ask what remains when the characteristic prime
is omitted. Our proposition supplies exactly that comparison.

For a finite field $k$ the group $J_X(k)$ is finite. The kernel
of its map to $\prod_{\ell\ne p}H^1(G_k,T_\ell J_X)$ is therefore
its $p$-primary subgroup. Equality of proper point sections
implies $[P-Q]\in J_X(k)[p^\infty]$. This is the principal-divisor
step used in A.19, with $P-Q$ obtained by subtracting the two
classes relative to the same base point; replacing it by $Q-P$
changes the sign but not its order.

More generally the proof of A.30 applies to \emph{semiabelian}
groups over fields finitely generated over their prime field.
We have rechecked it: after splitting the torus over a finite
separable extension, the divisor sequence makes $K^\times$
an extension of its constant units by a free abelian group;
the constants have finite torsion and a free abelian quotient.
The abelian-variety rational points are finitely generated by
Lang--N\'eron (\citealt[Theorem 2.1]{Conrad2006}) and the usual
Mordell--Weil theorem for number-field constants. Lifts of finitely
many generators and relations use only finitely many basis
elements of the free torus subgroup. Thus the semiabelian
rational-point group is a direct sum of a free abelian group
and a finitely generated group. Its intersection of all
prime-to-$p$ multiples is precisely its $p$-primary torsion,
and in characteristic zero the intersection of all multiples
is zero. Descent of the resulting torsion assertion through
the finite splitting extension is valid because rational-point
inclusion is injective. This proves the kernels asserted in A.30.

The claim on cote 149, scan 34, that all such $G(K)$ are finitely
generated is false for a torus: $\mathbf G_m(\mathbf Q)=\mathbf Q^\times$
has the independent valuations at all rational primes. The
divisibility conclusion above repairs the argument without that
claim. For a nonreduced modulus $J_{\mathfrak m}$ need not be
semiabelian; its unipotent part is invisible to prime-to-characteristic
Kummer theory. We therefore do not extend A.30 to that larger
class of groups. For example $J_{2\infty}$ of $\mathbf P^1$
is $\mathbf G_a$, on which every invertible $n$ acts bijectively.
